\documentclass[review]{elsarticle}

\usepackage{amsmath}
\usepackage{graphicx}
\usepackage{booktabs}
\usepackage{hyperref}

\journal{Target journal to be confirmed}

\begin{document}

\begin{frontmatter}

\title{How Robust Is Multitemporal Sentinel-1/Sentinel-2 Fusion for Flood Mapping under Optical Modality Degradation?}

\author[aff1]{Author Name}
\affiliation[aff1]{organization={Affiliation to be confirmed}}

\begin{abstract}
Multitemporal Sentinel-1/Sentinel-2 fusion can improve supervised flood mapping,
but disaster response often faces incomplete or degraded optical observations.
This study evaluates the sensitivity of flood segmentation models to sensor
modality availability using the public OMBRIA dataset. We compare Sentinel-1,
Sentinel-2, bitemporal, and multimodal U-Net-style baselines, then evaluate the
same multimodal checkpoint under controlled Sentinel-2 degradation at inference
time. The goal is not to claim a new state-of-the-art flood model, but to
quantify when fusion remains useful and when reliance on optical inputs becomes a
failure mode.
\end{abstract}

\begin{keyword}
flood mapping \sep Sentinel-1 \sep Sentinel-2 \sep multimodal fusion \sep optical degradation \sep remote sensing
\end{keyword}

\end{frontmatter}

\section{Introduction}

Draft after the first completed baseline experiments. The introduction should
motivate flood mapping, optical cloud limitations, SAR complementarity, and
modality robustness without claiming operational readiness.

\section{Related Work}

Cover OMBRIA and OmbriaNet \cite{drakonakis2022ombrianet}, cross-modal flood
mapping \cite{garg2023crossmodal}, and incomplete multimodal learning for remote
sensing \cite{chen2023incomplete}.

\section{Data}

Describe dataset version, GCS source, labels, bands, invalid masks, splits, and
event metadata.

\section{Methods}

Describe baselines, foundation-model adaptation, PEFT variants, training
protocol, and implementation details.

\section{Experimental Design}

Describe in-distribution, cross-event, cross-region, and low-label settings.
Define metrics before reporting results.

\section{Results}

Only insert metrics generated by committed configs and logs.

\section{Discussion}

Discuss practical value, failure modes, compute tradeoffs, and limitations.

\section{Conclusion}

Summarize evidence-backed findings.

\bibliographystyle{elsarticle-num}
\bibliography{../references/paper}

\end{document}
